\section{Related Work}
\label{sec:related_work}

The challenge of robust risk forecasting in digital psychiatry is inextricably linked to the problem of Informative Missingness (MNAR), where the mechanism of data loss is stochastically dependent on the latent clinical state of the patient \cite{nock2026predicting}. As digital phenotyping matures from feasibility studies to clinical decision support, the limitations of standard imputation and time-series modeling have become glaringly apparent. This section provides an exhaustive critique of the four dominant paradigms in high-dimensional clinical time-series imputation: adversarial distributional matching, recurrent smoothing, continuous-time differential equations, and score-based diffusion models. We evaluate these methods through the lens of the ``Digital Vacuum Paradox,'' where the predictive utility of a model is inversely proportional to its reliance on local temporal smoothness during periods of clinical disengagement.

\subsection{Adversarial Imputation and the Failure of Distributional Matching}

Generative Adversarial Networks (GANs) revolutionized missing data imputation by shifting the objective from point-wise error minimization to global distribution matching. The landmark GAIN framework \cite{yoon2018gain} introduced a hint-mechanism-based discriminator that forces the generator to produce imputed values $x_{imputed}$ such that the joint distribution $P(x_{observed}, x_{imputed})$ is indistinguishable from the underlying true distribution $P(X)$. Subsequent iterations, such as E2GAN \cite{luo2019e2gan}, integrated end-to-end discriminative loss with autoencoding reconstruction to improve stability.

However, in the context of psychiatric MNAR, adversarial matching suffers from a fundamental \textit{Alignment Collapse}. The GAN objective assumes that the observed data distribution is a representative, albeit incomplete, proxy for the ground truth. In clinical digital phenotyping, the missingness indicator $M_t$ is a downstream effect of the latent crisis state $S_t$. Because crises are statistically rare (the ``Imbalance Trap''), the discriminator learns to penalize any imputed values that deviate significantly from the ``healthy'' majority distribution. Consequently, when a patient stops engaging due to a depressive relapse, GAIN-based models tend to ``hallucinate'' the most likely state—usually a stable, non-crisis trajectory—thereby masking the very signal the clinician needs to detect. This failure is not merely a matter of hyperparameter tuning but a structural limitation of $f$-divergence minimization: in the absence of a causal prior $S \rightarrow M$, the model defaults to the mode of the observed distribution, effectively treating disengagement as a technical nuisance rather than a high-signal biomarker.

Furthermore, the adversarial framework relies on the assumption of a static data-generating process. In longitudinal psychiatric monitoring, the relationship between the latent state and the missingness mechanism can evolve (e.g., a patient may become more paranoid about tracking as their symptoms worsen). Current GAN-based approaches lack the longitudinal depth to capture this non-stationarity, leading to a ``Celerity Bias'' where the model overfits to the initial engagement patterns and fails to recognize the semantic shift in silence during acute episodes.

\subsection{Recurrent Models and the Smoothness Fallacy}

Prior to the rise of continuous-time models, Recurrent Neural Networks (RNNs) were the standard for handling clinical sequences. BRITS (Bidirectional Recurrent Imputation for Time Series) \cite{cao2018brits} and M-RNN \cite{yoon2018mrnn} utilize the hidden state $h_t$ to encode temporal dependencies, using bidirectional passes to ``anchor'' missing values between observed time points. BRITS, in particular, treats missing values as variables to be optimized during the backpropagation of the forecasting loss, effectively enforcing a local temporal consistency.

The failure of BRITS-like architectures in psychiatry stems from the \textit{Smoothness Fallacy}. Recurrent models, even those utilizing Gated Recurrent Units (GRUs) or Long Short-Term Memory (LSTM) cells, are biased toward continuous, gradual transitions in latent space. Psychiatric relapses, however, are often characterized by ``Jump'' transitions—acute, non-linear shifts in behavioral state (e.g., a sudden cessation of mobility or social interaction). When a large gap in data occurs (persistent disengagement), BRITS interpolates a smooth decay or growth between the last observed point and the next. This results in the ``Persistence Gap'': the model fails to capture the catastrophic state transition that occurred \textit{within} the vacuum. For instance, if a patient is stable on Monday and re-engages in a crisis on Friday, a recurrent model will suggest a linear decline throughout the week, whereas the reality might have been a stable plateau followed by a sharp cliff on Wednesday. By smoothing over these jumps, recurrent models provide a false sense of gradual clinical change, potentially delaying urgent intervention.

This local temporal dependency also makes recurrent models susceptible to ``The Local Trap.'' Because the hidden state is updated based on recent observations, a prolonged period of missing data causes the hidden state to either decay toward a default null vector or saturate, losing the causal context necessary for risk forecasting. Without a mechanism to distinguish between a ``meaningful pause'' and a ``technical failure,'' RNNs treat all gaps as equivalent, eroding the predictive utility of the informative missingness signal.

\subsection{Continuous-Time Neural ODEs/CDEs and Deterministic Limitations}

The introduction of Neural Ordinary Differential Equations (Neural ODEs) \cite{chen2018neuralode} and Neural Controlled Differential Equations (Neural CDEs) \cite{kidger2020neural} offered a mathematically rigorous approach to irregularly sampled time series. By modeling the latent state evolution as $dh/dt = f_{\theta}(h(t), t, x(t))$, these models allow for evaluation at any arbitrary time point, decoupling the model's resolution from the sampling frequency of the sensors. Li et al. \cite{li2020scalable} further extended this to Latent Stochastic Differential Equations (Latent SDEs), introducing a diffusion term to capture aleatoric uncertainty.

Despite their mathematical elegance, Neural CDEs are fundamentally limited by their \textit{Flow-Based Nature}. Standard Neural ODEs and CDEs assume that the latent trajectory is a continuous flow. While they can handle ``shocks'' via the driving signal $x(t)$, in the total absence of $x(t)$ (high-intensity missingness), the trajectory is governed entirely by the learned drift function $f_{\theta}$. In psychiatric applications, this leads to two critical issues. First, Neural CDEs are often too deterministic; without a robust jump-diffusion prior, the model cannot spontaneously ``jump'' to a high-risk state without an observed trigger. Second, as noted by Lu et al. \cite{lu2026neural_sdes_suicide_risk}, standard Neural ODEs lack a mechanism to expand uncertainty bands during disengagement.

While Latent SDEs address the uncertainty aspect, they still struggle with the \textit{Jump-Diffusion Discontinuity}: the diffusion term $\sigma(h, t)dW_t$ captures Brownian-like jitter but fails to account for the Poisson-distributed discrete jumps that represent acute relapse. Furthermore, the reliance on a continuous driving signal $X(t)$ (via natural cubic splines in CDEs) creates a ``hallucination artifact'' where the model must assume a path between observations to compute the integral. If the true path was a jump, the spline-based interpolation introduces a smoothing bias that underestimates the clinical volatility. Consequently, continuous-time models often produce ``over-smoothed'' risk trajectories that fail to satisfy the clinical requirement for detecting rapid, catastrophic disengagement.

\subsection{Score-based Diffusion Models and the Reversion to Healthy Noise}

The most recent frontier in clinical imputation (2024--2026) involves Conditional Score-based Diffusion Models, exemplified by CSDI \cite{tashiro2021csdi} and its clinical variants. These models treat imputation as a denoising process, learning the score function $\nabla_x \log p(x)$ to iteratively refine a noise-filled gap into a plausible signal. CSDI has shown remarkable performance in capturing complex multi-variate correlations in ICU data (e.g., MIMIC-III).

The specific failure mode of diffusion models in psychiatry is the \textit{Causal Priors Deficiency}. Diffusion models are inherently probabilistic and excellent at capturing the ``most plausible'' path. However, in psychiatric digital phenotyping, the ``most plausible'' path for the population is almost always the healthy path. Without an explicit causal model of the missingness mechanism ($M_t = f(S_t, \dots)$), the diffusion process lacks the ``clinical guidance'' necessary to move the imputed trajectory toward the high-risk region of the manifold during disengagement.

In 2024 and 2025, several attempts were made to integrate ``clinical prompts'' into the diffusion score function, yet these often resulted in ``hallucinated crises'' (Type I errors) or, more commonly, a reversion to ``healthy noise'' (Type II errors). The model effectively treats the missing gap as a region of maximum entropy and fills it with the mean of the training distribution. For a suicidal patient, this reversion to the mean is lethal. As argued by Cai et al. \cite{cai2026causal_estimands_n_of_1}, the missingness in digital phenotyping is not a corruption of the signal (noise) but an alternative representation of the signal itself. Current score-based models fail to recognize this equivalence, treating $M_t=1$ as a void to be filled rather than a measurement to be interpreted. This lack of causal grounding means that diffusion models prioritize distributional fidelity over clinical validity, which is unacceptable in high-stakes psychiatric decision-making.

\subsection{The Need for Causal Data Loss Prevention (DLP)}

The synthesis of these critiques reveals a systemic gap in the literature: a lack of models that treat missingness as a \textit{primary behavioral outcome}. Existing methods, from GAIN to CSDI, are reactive—they attempt to fix the data after it is lost. Our proposed Clinical DLP Framework shifts the paradigm to a proactive, causally grounded approach. By utilizing Jump-Diffusion priors and distinguishing between technical and clinical dropouts, we move beyond distribution matching and local smoothing.

We build upon the work of Joshi et al. \cite{joshi2021learning_to_defer} regarding Learning to Defer (L2D) and Parbhoo et al. \cite{parbhoo2022causal_ope} on Causal Off-Policy Evaluation, ensuring that when the sensors go silent, the model does not revert to a biased mean, but instead triggers a clinical vigilance protocol based on the predicted causal path of disengagement. Our framework addresses the ``Digital Exclusion'' risk by incorporating socio-technical covariates into the Clinical Vigilance Index (CVI), ensuring that patients with technical hardware limitations are not unfairly flagged as high clinical risk. This transition from ``imputation as recovery'' to ``missingness as inference'' is essential for the safety, equity, and reliability of ML-driven psychiatry in the 2026 clinical landscape.
